\documentclass[12pt,a4paper]{article}

% Русский язык
\usepackage[T2A]{fontenc}
\usepackage[utf8]{inputenc}
\usepackage[russian]{babel}

% Математика и графика
\usepackage{amsmath,amssymb}
\usepackage{graphicx}
\usepackage{float}

% Таблицы
\usepackage{tabularx}
\usepackage{booktabs}
\usepackage{array}
\usepackage{caption}
\captionsetup[table]{justification=raggedleft,singlelinecheck=false}
\setlength{\tabcolsep}{4pt}

% Ссылки и оглавление
\usepackage[hidelinks]{hyperref}
\usepackage{tocloft}

% Листинги кода
\usepackage{listings}
\usepackage{fancyvrb}
\lstset{
    basicstyle=\ttfamily\small,
    breaklines=true,
    frame=single,
    literate=
      {—}{{--}}1
      {←}{{$\leftarrow$}}1
      {→}{{$\rightarrow$}}1
}

% Геометрия и отступы
\usepackage{geometry}
\geometry{left=2.5cm,right=2.5cm,top=2cm,bottom=2cm}
\setlength{\parindent}{1.25cm}
\setlength{\parskip}{0pt}

% Формат заголовков
\usepackage{titlesec}
\titleformat{\section}{\large\bfseries}{\thesection}{1em}{}
\titleformat{\subsection}{\normalsize\bfseries}{\thesubsection}{1em}{}

\begin{document}

% Титульный лист
\begin{titlepage}
    \centering
    {\large \textbf{САНКТ-ПЕТЕРБУРГСКИЙ ГОСУДАРСТВЕННЫЙ УНИВЕРСИТЕТ}} \\[0.5cm]
    {\large Факультет прикладной математики--процессов управления} \\[0.5cm]
    {\large Программа бакалавриата} \\[0.2cm]
    {\large \textbf{«Большие данные и распределенная цифровая платформа»}} \\[4.5cm]

    {\Large \textbf{ОТЧЕТ}} \\[0.3cm]
    {\large По лабораторной работе № 2} \\[0.2cm]
    {\large По дисциплине «Алгоритмы и структуры данных»} \\[0.2cm]
    {\large На тему «Исследование задачи коммивояжера через алгоритм имитации отжига и муравьиный алгоритм с модификациями»} \\[4cm]

    \begin{flushright}
        Студент гр. 24Б15-пу \\
        Телицин А.Д. \\[0.5cm]
        Преподаватель \\
        Дик А.Г. \\[1cm]
    \end{flushright}

    \vfill
    \centering
    Санкт-Петербург \\
    2026 г.
\end{titlepage}

\tableofcontents
\newpage

\section{Цель работы}
\label{sec:goal}

Целью работы является исследование и сравнительный анализ качества решения задачи коммивояжера с помощью:
\begin{enumerate}
    \item базового алгоритма имитации отжига;
    \item модифицированного алгоритма имитации отжига (больцмановское охлаждение);
    \item базового муравьиного алгоритма;
    \item модифицированного муравьиного алгоритма («начальное расположение»).
\end{enumerate}

Основной акцент в сравнении делается на качестве получаемого маршрута (длине гамильтонова цикла), затем на времени выполнения и стабильности результатов при изменении параметров. Эксперименты проводятся на среднеразмерном графе Berlin52 и крупном графе World666.

\section{Описание алгоритмов и постановка задачи}
\label{sec:algorithms}

\subsection{Постановка задачи коммивояжера}

Пусть задан полный взвешенный граф $G=(V,E)$, где:
\begin{itemize}
    \item $V=\{1,2,\dots,n\}$ -- множество вершин (городов),
    \item $E$ -- множество ребер между всеми парами городов,
    \item $d_{ij}>0$ -- длина (вес) ребра между городами $i$ и $j$, $d_{ii}=0$.
\end{itemize}

Требуется найти гамильтонов цикл минимальной длины, то есть такую перестановку вершин
$\pi=(\pi_1,\pi_2,\dots,\pi_n)$, при которой минимизируется функционал:
\begin{equation}
L(\pi)=\sum_{k=1}^{n-1} d_{\pi_k,\pi_{k+1}} + d_{\pi_n,\pi_1} \rightarrow \min.
\end{equation}

В рамках данной лабораторной работы выполняется сравнение четырех алгоритмических вариантов:
\begin{enumerate}
    \item базовый алгоритм имитации отжига;
    \item модификация имитации отжига с больцмановским расписанием охлаждения;
    \item базовый муравьиный алгоритм;
    \item модификация муравьиного алгоритма «начальное расположение».
\end{enumerate}

\subsection{Алгоритм имитации отжига}

\textbf{Алгоритм имитации отжига} -- это стохастический метод комбинаторной оптимизации,
вдохновленный физическим процессом отжига металлов. В металлургии материал сначала
нагревают до высокой температуры, а затем постепенно охлаждают. При высоких температурах
атомы могут свободно переходить между состояниями, а при медленном охлаждении система
имеет шанс перейти в состояние с минимальной энергией.

В задаче коммивояжера проводится следующая аналогия:
\begin{itemize}
    \item состояние системы -- конкретный маршрут коммивояжера;
    \item энергия состояния -- длина маршрута;
    \item температура -- параметр, управляющий вероятностью принятия ухудшающих переходов.
\end{itemize}

Основная идея метода состоит в том, что на ранних этапах (при высокой температуре)
алгоритм имеет высокую степень свободы, чтобы не застревать в локальных
минимумах, а по мере охлаждения становится более статичным и стабилизируется около
лучших найденных решений.

Пусть:
\begin{itemize}
    \item $\pi$ -- текущее решение,
    \item $\pi'$ -- соседнее решение,
    \item $L(\pi)$, $L(\pi')$ -- длины маршрутов,
    \item $\Delta = L(\pi') - L(\pi)$,
    \item $T$ -- текущая температура.
\end{itemize}

Правило принятия решения:
\begin{equation}
\pi \leftarrow
\begin{cases}
\pi', & \Delta < 0, \\
\pi', & \Delta \ge 0 \text{ и } r < \exp\left(-\frac{\Delta}{T}\right), \\
\pi, & \text{иначе},
\end{cases}
\end{equation}
где $r \sim U(0,1)$.

Если новое решение улучшает целевую функцию ($\Delta < 0$), оно принимается всегда.
Если решение хуже ($\Delta \ge 0$), оно принимается с вероятностью Больцмана:
\begin{equation}
P(\text{accept}) = \exp\left(-\frac{\Delta}{T}\right).
\end{equation}

Для базовой версии используется геометрическое охлаждение:
\begin{equation}
T_{k+1} = \alpha T_k, \quad 0<\alpha<1.
\end{equation}

Для задачи коммивояжера критически важен способ генерации соседнего решения.
В данной работе используется следующий подход: выбираются две позиции маршрута, а
подмаршрут между ними разворачивается.

С практической точки зрения метод имитации отжига характеризуется следующими
особенностями:
\begin{itemize}
    \item хорошо подходит для сложных комбинаторных задач большой размерности;
    \item имеет небольшие вычислительные накладные расходы на одну итерацию;
    \item чувствителен к выбору расписания температуры и числу итераций;
    \item при слишком быстром охлаждении может рано зафиксироваться в локальном минимуме.
\end{itemize}

\subsection{Модификация имитации отжига: больцмановское охлаждение}

В модифицированном варианте используется другая функция температуры:
\begin{equation}
T(k) = \frac{T_0}{\ln(1+k)},
\end{equation}
где $k$ -- номер итерации ($k \ge 1$), $T_0$ -- начальная температура.

Вероятность принятия ухудшающего шага сохраняется:
\begin{equation}
P(\text{accept}) = \exp\left(-\frac{\Delta}{T(k)}\right), \quad \Delta>0.
\end{equation}

По сравнению с геометрическим охлаждением данная схема убывает существенно медленнее,
что повышает качество исследования пространства решений, но увеличивает время работы.
На практике это может привести к значтиельному улучшению качества решения, но при
этом требует больше времени на поиск.

\subsection{Муравьиный алгоритм}

\textbf{Муравьиный алгоритм} (Ant Colony Optimization, ACO) -- это эвристический
метод, основанный на идее коллективного поведения муравьев при поиске кратчайших
путей к источнику пищи.

В природной системе муравьи оставляют на пути химический след (феромон). Чем чаще
используется путь, тем больше на нем феромона, и тем вероятнее, что следующие муравьи
выберут именно его. При этом за счет испарения феромона система не «застывает» и
сохраняет способность к адаптации.

В задаче коммивояжера:
\begin{itemize}
    \item каждый муравей строит полный маршрут по графу;
    \item качество маршрута оценивается его длиной;
    \item более короткие маршруты усиливаются большим количеством феромона.
\end{itemize}

Обозначим:
\begin{itemize}
    \item $\tau_{ij}$ -- уровень феромона на ребре $(i,j)$;
    \item $\eta_{ij}=\frac{1}{d_{ij}}$ -- эвристическая привлекательность ребра;
    \item $\alpha$ -- степень влияния феромона;
    \item $\beta$ -- степень влияния эвристики;
    \item $\rho$ -- коэффициент испарения феромона;
    \item $Q$ -- коэффициент отложения феромона;
    \item $M$ -- число муравьев в колонии.
\end{itemize}

Вероятность перехода муравья из вершины $i$ в вершину $j$:
\begin{equation}
P_{ij}=
\frac{\tau_{ij}^{\alpha}\eta_{ij}^{\beta}}
{\sum\limits_{k \in \mathcal{N}_i} \tau_{ik}^{\alpha}\eta_{ik}^{\beta}},
\end{equation}
где $\mathcal{N}_i$ -- множество еще не посещенных вершин.

После завершения маршрутов всеми муравьями выполняется глобальное обновление феромона.
Сначала моделируется испарение:
\begin{equation}
\tau_{ij} \leftarrow (1-\rho)\tau_{ij}.
\end{equation}

Затем добавляется вклад всех муравьев:
\begin{equation}
\tau_{ij} \leftarrow \tau_{ij} + \sum_{m=1}^{M}\Delta\tau_{ij}^{(m)},
\end{equation}
где $M$ -- число муравьев, а вклад $m$-го муравья задается:
\begin{equation}
\Delta\tau_{ij}^{(m)}=
\begin{cases}
\frac{Q}{L_m}, & \text{если ребро }(i,j)\text{ входит в маршрут муравья }m,\\
0, & \text{иначе},
\end{cases}
\end{equation}
$L_m$ -- длина маршрута муравья $m$.

Итоговая формула обновления:
\begin{equation}
\tau_{ij} \leftarrow (1-\rho)\tau_{ij} + \sum_{m=1}^{M}\Delta\tau_{ij}^{(m)}.
\end{equation}

Таким образом, алгоритм сочетает два механизма:
\begin{itemize}
    \item Сохранение предыдущего опыта решений через отложение феромона, что способствует усилению хороших маршрутов.
    \item Способность к поиску новых решений за счет вероятностного выбора и испарения феромона, что предотвращает преждевременную сходимость.
\end{itemize}

\subsection{Модификация муравьиного алгоритма: «начальное расположение»}

В базовом варианте стартовая вершина каждого муравья выбирается случайно.
В модификации «начальное расположение» первая итерация организуется так, чтобы
муравьи стартовали из разных вершин более равномерно:
\begin{equation}
s_m = m \bmod n,
\end{equation}
где $s_m$ -- стартовая вершина муравья $m$, $n$ -- число вершин.

Далее алгоритм работает как базовый ACO.
Цель модификации -- улучшить начальное покрытие пространства решений и уменьшить
вероятность преждевременной концентрации феромона в ограниченной области графа.

\section{Пошаговое описание выполнения алгоритмов}
\label{sec:steps}

\subsection{Пошаговое описание: базовый алгоритм имитации отжига}

\begin{enumerate}
    \item \textbf{Инициализация.}
    Формируется начальный маршрут $\pi^{(0)}$, задаются начальная температура $T_0$,
    коэффициент охлаждения $\alpha$, максимальное число итераций $K_{\max}$.
    Вычисляется начальное значение целевой функции $L(\pi^{(0)})$.
    
    \item \textbf{Инициализация лучшего решения.}
    Текущее решение принимается как лучшее:
    $\pi^{best} \leftarrow \pi^{(0)}$, $L^{best} \leftarrow L(\pi^{(0)})$.

    \item \textbf{Генерация соседнего решения.}
    На итерации $k$ из текущего маршрута $\pi^{(k)}$ строится соседний маршрут
    $\pi'$ оператором 2-opt.

    \item \textbf{Оценка соседнего решения.}
    Вычисляется $\Delta = L(\pi') - L(\pi^{(k)})$.

    \item \textbf{Правило перехода.}
    Если $\Delta < 0$, переход выполняется всегда.
    Если $\Delta \ge 0$, переход выполняется с вероятностью
    $P=\exp\left(-\frac{\Delta}{T_k}\right)$.

    \item \textbf{Обновление лучшего решения.}
    Если новое текущее решение лучше глобально лучшего, то:
    $\pi^{best} \leftarrow \pi^{(k+1)}$, $L^{best} \leftarrow L(\pi^{(k+1)})$.

    \item \textbf{Охлаждение температуры.}
    Температура обновляется по формуле:
    $T_{k+1}=\alpha T_k$.

    \item \textbf{Проверка критерия остановки.}
    Если достигнуто $K_{\max}$ или $T_k < T_{\min}$, выполняется остановка.
    Иначе переход к шагу 3.

    \item \textbf{Завершение алгоритма.}
    Возвращается лучшее найденное решение $\pi^{best}$ и его длина $L^{best}$.
\end{enumerate}

\begin{figure}[H]
    \centering
    \includegraphics[width=0.9\textwidth]{images/sa_basic.drawio.png}
    \caption{Блок-схема базового алгоритма имитации отжига.}
    \label{fig:sa_basic_flow}
\end{figure}

\subsection{Пошаговое описание: имитация отжига с больцмановским охлаждением}

\begin{enumerate}
    \item \textbf{Инициализация.}
    Формируется начальный маршрут $\pi^{(0)}$, задаются $T_0$ и $K_{\max}$.

    \item \textbf{Инициализация лучшего решения.}
    Устанавливается:
    $\pi^{best} \leftarrow \pi^{(0)}$, $L^{best} \leftarrow L(\pi^{(0)})$.

    \item \textbf{Вычисление температуры по номеру итерации.}
    Для текущей итерации $k$:
    \[
    T(k)=\frac{T_0}{\ln(1+k)}.
    \]

    \item \textbf{Генерация соседнего решения.}
    Строится $\pi'$ из $\pi^{(k)}$ оператором 2-opt.

    \item \textbf{Оценка перехода и принятие решения.}
    Вычисляется $\Delta = L(\pi') - L(\pi^{(k)})$.
    Применяется то же правило принятия, что и в базовом SA.

    \item \textbf{Обновление лучшего найденного решения.}
    При улучшении обновляются $\pi^{best}$ и $L^{best}$.

    \item \textbf{Проверка критерия остановки.}
    Если $k \ge K_{\max}$ или $T(k) < T_{\min}$, переход к завершению.
    Иначе переход к шагу 3 следующей итерации.

    \item \textbf{Завершение алгоритма.}
    Возвращается $\pi^{best}$ и $L^{best}$.
\end{enumerate}

\begin{figure}[H]
    \centering
    \includegraphics[width=0.9\textwidth]{images/sa_boltzmann.drawio.png}
    \caption{Блок-схема алгоритма имитации отжига с больцмановским охлаждением.}
    \label{fig:sa_boltzmann_flow}
\end{figure}

\subsection{Пошаговое описание: базовый муравьиный алгоритм}

\begin{enumerate}
    \item \textbf{Инициализация параметров.}
    Задаются $M$ (число муравьев), $I_{\max}$ (число итераций),
    $\alpha$, $\beta$, $\rho$, $Q$.

    \item \textbf{Инициализация феромона.}
    Для всех ребер $(i,j)$ задается начальный уровень феромона $\tau_{ij}^{(0)}$.
    Вычисляется эвристика $\eta_{ij}=1/d_{ij}$.

    \item \textbf{Начало итерации колонии.}
    Для итерации $t$ каждый муравей начинает построение маршрута
    из стартовой вершины (в базовой версии старт выбирается случайно).

    \item \textbf{Построение маршрутов муравьями.}
    Каждый муравей пошагово выбирает следующую непосещенную вершину
    по вероятностному правилу $P_{ij}$ до построения полного гамильтонова цикла.

    \item \textbf{Оценка маршрутов.}
    Для каждого муравья вычисляется длина маршрута $L_m$.
    Обновляется глобально лучший маршрут при необходимости.

    \item \textbf{Испарение феромона.}
    Выполняется:
    \[
    \tau_{ij}\leftarrow(1-\rho)\tau_{ij}.
    \]

    \item \textbf{Отложение феромона.}
    Для всех муравьев добавляется вклад:
    \[
    \tau_{ij}\leftarrow\tau_{ij}+\sum_{m=1}^{M}\Delta\tau_{ij}^{(m)}.
    \]

    \item \textbf{Проверка критерия остановки.}
    Если $t \ge I_{\max}$, переход к завершению.
    Иначе переход к шагу 3.

    \item \textbf{Завершение алгоритма.}
    Возвращается лучший найденный маршрут и его длина.
\end{enumerate}

\begin{figure}[H]
    \centering
    \includegraphics[width=0.9\textwidth]{images/aco_base.drawio.png}
    \caption{Блок-схема базового муравьиного алгоритма.}
    \label{fig:aco_basic_flow}
\end{figure}

\subsection{Пошаговое описание: муравьиный алгоритм с модификацией «начальное расположение»}

\begin{enumerate}
    \item \textbf{Инициализация параметров и феромона.}
    Выполняется аналогично базовому ACO.

    \item \textbf{Специальная инициализация стартовых вершин (первая итерация).}
    На первой итерации старт муравья $m$ задается по правилу:
    \[
    s_m = m \bmod n.
    \]
    Это обеспечивает равномерное покрытие графа в начале работы алгоритма.

    \item \textbf{Построение маршрутов и их оценка.}
    Выполняется как в базовом ACO с обновлением лучшего решения.

    \item \textbf{Обновление феромона.}
    Испарение и отложение выполняются как в базовом ACO.

    \item \textbf{Последующие итерации.}
    Начиная со второй итерации алгоритм работает как базовый ACO
    (со случайным или стандартным стартом, в зависимости от реализации).

    \item \textbf{Проверка критерия остановки.}
    Если достигнуто $I_{\max}$, переход к завершению.
    Иначе повторение цикла итераций.

    \item \textbf{Завершение алгоритма.}
    Возвращается лучший найденный маршрут и его длина.
\end{enumerate}

\begin{figure}[H]
    \centering
    \includegraphics[width=0.9\textwidth]{images/aco_init.drawio.png}
    \caption{Блок-схема муравьиного алгоритма с модификацией «начальное расположение».}
    \label{fig:aco_initial_flow}
\end{figure}

\section{Рекомендации пользователю}
\label{sec:user_recommendations}

Взаимодействие с программой осуществляется через графический интерфейс.

\begin{enumerate}
    \item \textbf{Запуск программы.} Активируйте виртуальное окружение и выполните команду:
    \begin{lstlisting}[language=bash]
python3 main.py
    \end{lstlisting}

    \item Откроется окно, где можно выбрать граф, алгоритм и параметры.
\begin{figure}[H]
    \centering
    \includegraphics[width=0.9\textwidth]{images/ui_1.png}
    \caption{Интерфейс программы.}
    \label{fig:aco_initial_flow}
\end{figure}

    \item \textbf{Выбор алгоритма.} Производится через выпадающий список в интерфейсе.
    На выбор представлены четыре варианта: базовый и больцмановский отжиг, базовый и с
    начальным расположением муравьиный алгоритм.

    \item \textbf{Выбор графа.} Производится через выпадающий список. Доступные графы: \textit{small}, \textit{berlin52} и \textit{world666}.

\begin{figure}[H]
    \centering
    \includegraphics[width=0.9\textwidth]{images/ui_2.png}
    \caption{Интерфейс программы.}
    \label{fig:aco_initial_flow}
\end{figure}

    \item \textbf{Настройка параметров.} Программа позволяет изменять ключевые параметры алгоритмов,
    такие как начальная температура, коэффициент охлаждения для отжига, а также $\alpha$, $\beta$, $\rho$
    и $Q$ для муравьиного алгоритма. Рекомендуется начинать с базовых значений и постепенно их изменять.

    \item \textbf{Результаты.} После запуска следует ориентироваться не только
    на длину найденного маршрута, но и на график сходимости. Если кривая быстро
    выходит на плато, это означает, что выбранные параметры слишком быстро
    ограничивают поиск.

\begin{figure}[H]
    \centering
    \includegraphics[width=0.9\textwidth]{images/ui_3.png}
    \caption{Интерфейс программы.}
    \label{fig:aco_initial_flow}
\end{figure}

\end{enumerate}

\section{Рекомендации программисту}
\label{sec:developer_recommendations}

\textbf{Системные требования:}
\begin{itemize}
    \item операционная система macOS, Linux или Windows;
    \item Python 3.14 или совместимая версия Python 3;
    \item установленный пакет PyQt6 для запуска графического интерфейса;
    \item установленный NumPy для вычислений и построения раскладки графа.
\end{itemize}

\textbf{Структура проекта:}
\begin{itemize}
    \item \texttt{src/graph\_utils.py} -- загрузка графов и вычисление длины маршрута;
    \item \texttt{src/annealing.py} -- реализация базового и больцмановского отжига;
    \item \texttt{src/aco.py} -- реализация базового ACO и модификации initial placement;
    \item \texttt{src/ui.py} -- графический интерфейс, запуск алгоритмов и визуализация;
    \item \texttt{main.py} -- точка входа в приложение.
\end{itemize}


\textbf{Инструкция по запуску:}
\begin{enumerate}
    \item Активировать виртуальное окружение.
    \item Установить зависимости из \texttt{requirements.txt}.
    \item Запустить \texttt{python3 main.py}.
    \item Выбрать граф, алгоритм и параметры в окне программы.
\end{enumerate}

\section{Описание контрольного примера}
\label{sec:control_example}

В качестве контрольного примера использован малый граф \textit{small} из шести
вершин. Этот пример удобен тем, что он не требует значительных вычислительных ресурсов,
а оптимальный маршрут известен заранее (длина 10). Это позволяет проверить
корректность реализации алгоритмов и убедиться,

В ходе проверки были получены следующие результаты:

\begin{table}[H]
    \centering
    \caption{Результаты контрольного прогона на графе \textit{small}.}
    \begin{tabularx}{\textwidth}{p{5cm}rrr}
        \toprule
        Алгоритм & Длина маршрута & Итерации & Время, сек \\
        \midrule
        Базовая имитация отжига & 10 & 225 & 0.0023 \\
        Имитация отжига с больцмановским охлаждением & 10 & 1000 & 0.0093 \\
        Базовый муравьиный алгоритм & 10 & 50 & 0.0419 \\
        Муравьиный алгоритм + начальное расположение & 10 & 50 & 0.0669 \\
        \bottomrule
    \end{tabularx}
\end{table}

\begin{figure}[H]
    \centering
    \includegraphics[width=0.9\textwidth]{images/ui_4.png}
    \caption{Интерфейс программы.}
    \label{fig:aco_initial_flow}
\end{figure}


Полученные значения подтверждают корректность реализации: все четыре варианта
алгоритмов находят оптимальный маршрут длины 10 на контрольном графе.
Визуализация в графическом интерфейсе показывает найденный маршрут на двумерной
раскладке графа. График сходимости найденного решения демонстрирует, что алгоритмы
повышают качество решения по мере итераций.

\section{Исследование параметров и анализ результатов}
\label{sec:research}

\subsection{Методика исследования}

Для оценки влияния параметров на качество решения были выполнены серии экспериментов
для графов Berlin52 и World666. В каждой серии изменялся один параметр, а остальные
оставались фиксированными. В качестве метрик, на основе которых проводилось сравнение
алгоритмов были выбраны следующие (в порядке убывания значимости):
\begin{itemize}
    \item \textbf{Длина найденного маршрута} -- основная метрика качества решения.
    \item \textbf{Время работы} -- важный показатель эффективности алгоритма.
    \item \textbf{Устойчивость результатов} -- оценка вариативности решения при повторных запусках с разными начальными условиями.
\end{itemize}

Для Berlin52 эксперименты выполнялись с небольшим числом повторов, а для World666
число повторов было сокращено, в целях снижения общего времени вычислений. Поэтому
при интерпретации результатов основное внимание уделяется не абсолютной точности
статистики, а общим тенденциям поведения алгоритмов.

\subsection{Имитация отжига (на Berlin52)}

В базовом варианте имитации отжига наибольшее влияние на качество решения оказал
выбор начальной температуры $T_0$ и коэффициента охлаждения $\alpha$.
На Berlin52 оптимальные результаты были получены при $T_0=800$ и $\alpha=0.9989$,
что обеспечивало достаточно медленное охлаждение. При более быстром охлаждении
($\alpha < 0.95$) в сочетании с низкой начальной температурой алгоритм часто застревал
в локальных минимумах. Полученное расстояние составило 1980, но при этом данный результат
не был стабильным и при повторных запусках мог значительно ухудшаться. Время выполнения
при данных параметрах составило 0.1893 секунды, что является приемлемым для данной задачи.

В варианте с больцмановским охлаждением оптимальные результаты были достигнуты при
$T_0=40$ и $K_{\max}=10000$. Этот вариант показал немного более стабильные результаты,
с меньшей вариативностью между запусками, лучшее найденное расстояние составило 1941,
а время выполнения осталось приблизительно тем же, что и в варианте без модификации,
0.1832 секунды. Особенностью больцмановского охлаждения является более медленное
снижение температуры, что позволяет алгоритму более тщательно исследовать пространство
решений, но при этом требует другого подхода к выбору параметров, так например,
$T_0=40$ является низким относительно базового отжига, но при этом обеспечивает более
эффективный поиск при медленном охлаждении. Более высокие значения $T_0$ в данном случае
приводили к избыточной случайности и ухудшению результатов.

\subsection{Имитация отжига (на World666)}

Базовый алгоритм имитации отжига на World666 показал более высокую чувствительность к параметрам.
Наилучшие результаты были получены при $T_0=800$ и $\alpha=0.9999$, что обеспечивало достаточно
медленное охлаждение для такого большого графа. Полученное расстояние составило 584576, а время
выполнения -- 16.0894 секунды. Число ограничения итераций было выбрано таким образом, чтобы
алгоритм мог стабилизироваться около лучших решений, что особенно важно при медленном охлаждении.
При более быстром охлаждении алгоритм демонстрировал более высокую вариативность и хуже качество
решения, а при более низких $T_0$ он часто застревал в локальных минимумах.

Алгоритм в варианте с больцмановским охлаждением показал наиболее качественные решения при
$T_0=60$ и $K_{\max}=100000$. Полученное расстояние составило 609483. Время выполнения
составило 15.8380 секунды. Выбор таким параметров был обусловлен размером графа и
необходимостью обеспечить достаточно медленное охлаждение для эффективного поиска.
При более высоких $T_0$ алгоритм демонстрировал более высокую вариативность и хуже
качество решения, а при более низких $T_0$ он часто застревал в локальных минимумах.
Ограничение итераций было выбрано таким образом, чтобы алгоритм мог стабилизироваться
около лучших решений, что особенно важно при медленном охлаждении.

Таким образом, для имитации отжига можно сделать следующие выводы:
\begin{itemize}
    \item оба алгоритма позволяют достичь хороших результатов при правильной настройке параметров, однако
    \item наибольшее влияние на характер работы алгоритма оказывает выбор и настройка функции температуры;
    \item слишком быстрое охлаждение ухудшает поиск на обоих графах;
    \item больцмановская схема требует больше итераций для достижения хорошего результата, но при этом может значительно улучшить качество решения;
    \item больцмановская схема требует иного подхода к выбору начальной температуры, которая должна быть достаточно низкой для эффективного поиска при медленном охлаждении;
    \item важным является выбор ограничений итераций таким образом, чтобы алгоритм
работал в том числе и при более низких температурах, что позволяет ему стабилизироваться около лучших решений.
\end{itemize}

\subsection{Муравьиный алгоритм (на Berlin52)}

Для муравьиного алгоритма наиболее значимым оказался баланс между влиянием феромона
и эвристики. На обеих задачах улучшение результатов наблюдалось при увеличении
$\alpha$ и $\beta$ по сравнению с базовыми значениями. Это означает, что алгоритму
полезно сильнее учитывать как уже найденные хорошие ребра, так и текущую привлекательность
коротких переходов.

На Berlin52 наилучшие результаты среди проверенных серий были получены при количестве
муравьев $M=30$, количестве итераций $I_{\max}=50$,
$\alpha = 1.0$ и $\beta = 2.0$. При этом уменьшение или увеличение $\rho$ вокруг
$0.3-0.6$ приводило к ухудшению результата, то есть слишком слабое испарение приводило к
ранней фиксации феромона, а слишком сильное -- к потере накопленной информации. Полученный
результат составил 1928, а время выполнения -- 1.3808 секунды.

С использованием модификации «начальное расположение» на Berlin52 были получены лучшие
результаты при $M=30$, $I_{\max}=50$, $\alpha = 1.5$, $\beta = 2.0$ и $Q = 50$.
Этот вариант показал улучшение качества решения, что может быть связано с более
равномерным начальным покрытием графа, что позволяет алгоритму эффективнее исследовать
пространство решений. Полученный результат составил 1942,
а время выполнения -- 1.3739 секунды.

\subsection{Муравьиный алгоритм (на World666)}

На World666 испытания были осложнены большим размером графа, что ограничивало число повторов.
Тем не менее, были подобраны следующие параметры: $M=30$, $I_{\max}=50$, $\alpha = 1.5$,
$\beta = 2.0$ и $\rho = 0.4$, что является той же самой конфигурацией, что и для Berlin52.
Полученное расстояние составило 397880, а время выполнения -- 131.4481 секунды.

Для модификации «начальное расположение» характерна более равномерная начальная
разведка графа. На Berlin52 это дало небольшое преимущество на некоторых параметрах,
а на World666 различия между базовым и модифицированным вариантом были менее
однозначными. Тем не менее, модификация полезна как способ уменьшить зависимость
результата от случайного старта. Полученный результат для World666 составил 401051,
а время выполнения -- 171.5460 секунды при использовании параметров $M=30$,
$I_{\max}=50$, $\alpha = 1.5$, $\beta = 3.0$, $\rho = 0.4$ и $Q = 100$.

\subsection{Сводная таблица результатов}

\begin{table}[H]
    \centering
    \caption{Ключевые результаты исследования параметров на графе Berlin52.}
    \begin{tabularx}{\textwidth}{p{3.6cm}Xr}
        \toprule
        Алгоритм & Наиболее удачные параметры в серии испытаний & Длина маршрута \\
        \midrule
        Имитация отжига (базовый) & $\alpha = 0.9989$; $T_0 = 800$ & 1980 \\
        Имитация отжига (Больцман) & $T_0 = 40$; $K_{\max} = 10000$ & 1941 \\
        Муравьиный алгоритм (базовый) & $\alpha = 1.0$; $\beta = 2.0$; $\rho = 0.4$ & 1928 \\
        Муравьиный алгоритм (начальное расположение) & $\alpha = 1.5$; $\beta = 2.0$; $Q = 50$ & 1942 \\
        \bottomrule
    \end{tabularx}
\end{table}

\begin{table}[H]
    \centering
    \caption{Ключевые результаты исследования параметров на графе World666.}
    \begin{tabularx}{\textwidth}{p{3.6cm}Xr}
        \toprule
        Алгоритм & Наиболее удачные параметры в серии испытаний & Длина маршрута \\
        \midrule
        Имитация отжига (базовый) & $\alpha = 0.9999$; $T_0 = 800$ & 584576 \\
        Имитация отжига (Больцман) & $T_0 = 60$; $K_{\max} = 100000$ & 609483 \\
        Муравьиный алгоритм (базовый) & $\alpha = 1.5$; $\beta = 2.0$; $\rho = 0.4$ & 397880 \\
        Муравьиный алгоритм (начальное расположение) & $\alpha = 1.5$; $\beta = 3.0$; $\rho = 0.4$; $Q = 100$ & 401051 \\
        \bottomrule
    \end{tabularx}
\end{table}

\subsection{Сравнение алгоритмов}

С точки зрения \textbf{качества} на графе Berlin52 оба алгоритма обеспечили
поиск решения, близкого к оптимальному (1928 против 1941-1942), при этом
муравьиный алгоритм показал чуть лучшее качество при больших временных затратах.

На графе World666 разница в качестве стала более заметной:
муравьиный алгоритм обеспечил значительно более короткий маршрут
(397880 против 584576-609483), что свидетельствует о его лучшей
способности находить хорошие решения на больших и сложных графах. При этом
стоит учитывать, что время выполнения муравьиного алгоритма было существенно
больше, что может быть критичным в условиях ограниченных ресурсов.

С точки зрения \textbf{скорости} имитация отжига работает существенно быстрее, чем
муравьиный алгоритм, поскольку на каждой итерации обрабатывает только одно соседнее
решение, тогда как ACO строит множество маршрутов сразу для всей колонии.

С точки зрения \textbf{стабильности} муравьиный алгоритм показал более устойчивые
результаты при повторных запусках, особенно на большом графе, что может быть
связано с его механизмом коллективного обучения через феромоны. Еще сильнее увеличивает
стабильность модификация «начальное расположение», которая обеспечивает более
равномерное начальное покрытие графа.

Имитация отжига, особенно в базовом варианте, была более чувствительна
к начальным условиям и параметрам, что приводило к большей вариативности
результатов. Модификация с больцмановским охлаждением улучшила стабильность,
но при этом потребовала более тщательного подбора параметров и увеличила время выполнения.

В результате можно сделать итоговый вывод: допустимо использование обоих алгоритмов,
однако при приоритете качества решения для задач средней и большой размерности
предпочтителен муравьиный алгоритм с тщательным подбором параметров.
При приоритете времени вычислений может использоваться имитация отжига,
но это обычно приводит к ухудшению качества решения на сложных графах.

\section{Выводы}

В ходе лабораторной работы была исследована задача коммивояжера с применением
четырех вариантов алгоритмов: базовой имитации отжига, имитации отжига с
больцмановским охлаждением, базового муравьиного алгоритма и муравьиного
алгоритма с модификацией «начальное расположение».

По результатам экспериментов на графах Berlin52 и World666 получены следующие
основные выводы.

\begin{enumerate}
    \item \textbf{По критерию качества решения} на обеих задачах муравьиный алгоритм
    показал лучшие результаты, чем имитация отжига. Разница особенно заметна на
    графе World666, где ACO дает существенно более короткие маршруты.

    \item \textbf{По критерию скорости} имитация отжига работает быстрее,
    поскольку на каждой итерации оценивается одно соседнее решение, тогда как
    муравьиный алгоритм строит набор маршрутов всей колонией.

    \item \textbf{По критерию стабильности} и управляемости результата ключевую роль
    играют параметры алгоритмов. Для имитации отжига наиболее критичен закон
    охлаждения (в первую очередь коэффициент охлаждения в базовой версии). Для ACO
    наиболее значимы параметры $\alpha$ и $\beta$, задающие баланс между феромоном и
    эвристической информацией.

    \item Больцмановская модификация отжига обеспечивает более глубокий поиск
    за счет медленного снижения температуры и на сложных графах может заметно
    улучшать качество по сравнению с базовой версией, но ценой увеличения времени.

    \item Модификация муравьиного алгоритма «начальное расположение» улучшает
    начальное покрытие графа и в ряде режимов дает преимущество по качеству
    решения, особенно на средних размерах задачи.
\end{enumerate}

Итоговый практический вывод: при приоритете \textbf{качества решения} для задачи
коммивояжера на графах средней и большой размерности предпочтителен муравьиный
подход с подбором параметров. При приоритете \textbf{времени вычислений} может
использоваться имитация отжига, однако это обычно приводит к ухудшению качества
маршрута на сложных графах.

Полученные результаты корректны для конкретных реализаций алгоритмов и выбранных графов,
и могут отличаться при других настройках или на других задачах.
Тем не менее, общие тенденции в поведении алгоритмов сохраняются, что позволяет сделать
обоснованные рекомендации по их применению в зависимости от требований к качеству
решения и доступного времени.

\pagebreak

\section{Листинг кода}
В данном разделе приведены исходные коды основных модулей программы.

\subsection*{Листинг 1. Исходный код файла \texttt{main.py}}
\VerbatimInput[fontsize=\small, frame=single]{report_code/main.py}

\subsection*{Листинг 2. Исходный код файла \texttt{src/graph\_utils.py}}
\VerbatimInput[fontsize=\small, frame=single]{report_code/src_graph_utils.py}

\subsection*{Листинг 3. Исходный код файла \texttt{src/annealing.py}}
\VerbatimInput[fontsize=\small, frame=single]{report_code/src_annealing.py}

\subsection*{Листинг 4. Исходный код файла \texttt{src/aco.py}}
\VerbatimInput[fontsize=\small, frame=single]{report_code/src_aco.py}

\subsection*{Листинг 5. Исходный код файла \texttt{src/ui.py}}
\VerbatimInput[fontsize=\small, frame=single]{report_code/src_ui.py}

\end{document}
